\section{Context: fitting 262{,}144 tokens into two 16\,GB budgets}\label{sec:context}

\subsection{What the KV cache costs on this architecture}
Qwen3.8-27B is a hybrid model: of its 64 layers only every fourth uses full attention, and only those
16 layers keep a KV cache. With 4 KV heads of dimension 256, a token costs $2\times16\times4\times256
= 32{,}768$ cached elements. In llama.cpp's block formats that is exactly \textbf{64\,KiB per token at
\flag{f16}, 34\,KiB at \flag{q8\_0} and 18\,KiB at \flag{q4\_0}}, summed over both cards. At the native
262{,}144 tokens the cache is 16.0, 8.5 or 4.5\,GiB. The server log confirms it
(\texttt{llama\_kv\_cache: size = 4608.00 MiB (262144 cells, 16 layers)} at \flag{q4\_0}). The
remaining 48 layers carry a fixed recurrent state of 449\,MiB regardless of length. MTP adds a
one-layer draft cache that is always \flag{f16}: 1{,}024\,MiB at the full window. A dense 27B model
with 64 attention layers would need four times as much, so the hybrid design is what makes 262K tokens
thinkable on this hardware.

An earlier ``VRAM model'' in my own notes, fitted to measurements from the deleted tensor-split build,
had per-GPU rates that were not per-token KV at all, and predicted ceilings that did not reproduce. The
arithmetic above replaced it. It is only part of the footprint: compute buffers, the drafter and
per-card placement decide the real ceiling, which has to be measured.

\subsection{The default split strands memory; a manual split recovers it}\label{sec:ts}

\begin{figure}[t]
\centering
\begin{tikzpicture}
\begin{axis}[paper,xbar stacked,height=5.0cm,width=0.66\linewidth,bar width=9pt,
  ytick={1,2,4,5},ymin=0.3,ymax=5.7,
  yticklabels={{default split (ceiling 196{,}608) \ GPU1},{GPU0},{\flag{-ts 58,42} (ceiling 262{,}144) \ GPU1},{GPU0}},
  yticklabel style={font=\scriptsize},
  xmin=0,xmax=17500,xlabel={peak VRAM after deep prefill (MiB)},
  scaled x ticks=false,xticklabel style={/pgf/number format/1000 sep={,}},
  legend style={at={(0.5,1.04)},anchor=south,legend columns=2}]
\addplot[fill=cbblue!75,draw=cbblue!60!black] coordinates {(15640,1) (12978,2) (15082,4) (15322,5)};
\addplot[fill=black!10,draw=black!40,postaction={pattern=north east lines,pattern color=black!35}] coordinates {(671,1) (3333,2) (1229,4) (989,5)};
\legend{used,free (unreachable from the other card)}
\draw[cbred,dashed,line width=0.8pt] (axis cs:16311,0.3) -- (axis cs:16311,5.7);
\end{axis}
\end{tikzpicture}
\caption{\Qsix{} with MTP and \flag{q4\_0} KV, build E-A, greedy. At the engine's default layer split
the largest window that worked was 196{,}608 tokens: GPU1 was 671\,MiB from its limit while 3.3\,GiB
sat idle on GPU0. Moving layers to GPU0 with \flag{-ts 58,42} reaches the native 262{,}144 with both
cards within 1.3\,GiB of full. The dashed line is the 16{,}311\,MiB card limit.}\label{fig:vram}
\end{figure}

Under layer split the last card also holds the output head and, with MTP, the draft context, so the
default placement overloads GPU1. \Cref{fig:vram} shows the effect for \Qsix. One flag moved the ceiling
from 196{,}608 to 262{,}144 tokens.

\begin{table}[t]
\centering\small
\caption{Context ceilings, first battery: MTP $n{=}2$, \flag{q4\_0} KV, layer split, \flag{-fit off},
official non-thinking sampling, 192-token generation after a $\geq$94.7\,\% prefill. Failures at
262{,}144 are single attempts except where noted; all are compute-buffer out-of-memory errors read from
the saved server log.}\label{tab:ceilings}
\shrinktofit{\begin{tabular}{@{}lrlllr@{}}
\toprule
build & ceiling & splits that load & splits that failed & engine default & GPU0 / GPU1 MiB \\
\midrule
\QsixXL & 212{,}992\textsuperscript{a} & 56,44 & (at 196K) default, 52,48, 54,46, 58,42 & not tried at ceiling & 15{,}656 / 15{,}414 \\
\Qsix   & 262{,}144 & 58,42; 56,44 & default, 54,46, 60,40, 62,38 & \textbf{fails} & 15{,}322 / 15{,}082 \\
\Qfive  & 262{,}144 & 54,46 to 62,38 (five) & default & \textbf{fails} & 14{,}660 / 15{,}402 \\
\Qfour  & 262{,}144 & default; 56,44 & --- & \textbf{loads} & 12{,}116 / 15{,}094 \\
\bottomrule
\multicolumn{6}{@{}l}{\footnotesize\textsuperscript{a}\,Bracketed: 229{,}376 failed twice. The other three reach the model's maximum, so no rung exists above them.}
\end{tabular}}
\end{table}

\Cref{tab:ceilings} gives all four builds. Three findings.

\textbf{For three of four builds the split, not the quantization, set the reachable window.} \Qfive{}
and \Qsix{} fail at the default split and reach the native maximum once rebalanced; \QsixXL{} goes from
131{,}072 (its earlier default-split ceiling) to 212{,}992. The smallest build is the counter-example:
it loads at the default. I first wrote this finding as ``the ceiling is a property of the split'' for
the whole ladder; review found the counter-example in my own data. A published context ceiling without
its split is incomplete either way.

\textbf{The best split does not transfer.} \Qsix{} wants \flag{58,42}. On \QsixXL{} that same ratio
overshoots (GPU0-heavy by 2.2\,GiB) and only \flag{56,44} loads; on \Qsix, \flag{54,46} fails while
ratios on both sides of the working pair also fail. The relationship is not monotone. Re-sweep after
any change of build, KV precision or speculative setting.

\textbf{Balanced is not fastest.} On \Qfive{} at the full window the most balanced split (28\,MiB
apart) decoded at 8.50\tps; four other splits with 166 to 1{,}750\,MiB of imbalance clustered at
10.4--10.8\tps. The outlier is a single reading and unexplained. The rule that survives: keep the
fastest split that loads, not the most balanced one. Otherwise the split has little effect on speed:
with the generated text held fixed, it moved decode by at most 5.7\,\% in 11 of 12 groups.

\subsection{The full map: context $\times$ KV precision $\times$ drafter}\label{sec:kvmap}

\begin{table}[t]
\centering\small
\caption{Largest window that loads, reports the requested size, prefills $\geq$94\,\% and generates
512 tokens. Second battery, upstream build E-C, layer split at each cell's best ratio, 71 launch
attempts; every ceiling below the native window is bracketed by a rung above it that failed twice. DFlash2 cells place the drafter
on GPU1 (\flag{-devd CUDA1}). Bold: the full native window.}\label{tab:kvmap}
\shrinktofit{\begin{tabular}{@{}llrrr@{}}
\toprule
build & drafter & \flag{q4\_0} KV (18 KiB/tok) & \flag{q8\_0} KV (34 KiB/tok) & \flag{f16} KV (64 KiB/tok) \\
\midrule
\Qsix & MTP & \textbf{262{,}144} & 196{,}608 & 131{,}072 \\
\Qsix & DFlash2 & \textbf{262{,}144} & 196{,}608 & 98{,}304 \\
\Qfour & MTP & \textbf{262{,}144} & \textbf{262{,}144} & 196{,}608 \\
\Qfour & DFlash2 & \textbf{262{,}144} & \textbf{262{,}144} & 163{,}840 \\
\bottomrule
\end{tabular}}
\end{table}

\Cref{tab:kvmap} is the configuration space the rest of the second battery draws from. Two
consequences. The 4.1\,GiB saved by \Qfour{} over \Qsix{} buys a whole step of KV precision at the
same window: \Qfour{} holds the native 262K with an 8-bit cache, \Qsix{} cannot. Given
\cref{sec:kvq}, that is a real trade: weight fidelity against cache fidelity. And the two drafters
reach the same window except with an unquantized cache, where DFlash2's separate 1.1\,GB model costs
one rung. Two passing cells peaked above 15{,}700\,MiB on a card, past what I had treated as a
practical limit; the soak in \cref{sec:hostram} is what tests whether such cells are stable.

\subsection{Three ways a long-context server fails}
A ceiling search has to distinguish at least three failure modes, all seen on this host:
\begin{enumerate}
\item \textbf{Out of memory at load.} The clean case: the server log names the buffer and the device.
\item \textbf{Hang during initialisation.} \Qfour{} at 229K and 262K on the default split, in the
ceiling runs that preceded the first battery: no error in ten minutes, a 572-byte log. Not an OOM, and not to be labelled
one.
\item \textbf{Death at the first speculative draft.} \Qfour{} with DFlash2 and \flag{f16} KV at
196{,}608 tokens: the server loaded, prefilled all 184{,}807 tokens at 612\tps{}, then aborted inside
the drafter on a cuBLAS workspace allocation with 15.8\,GiB in use on each card. Twice. The client saw
only a dropped connection. A search that checks loading alone, or reads a disconnect as a client
fault, records this rung as usable.
\end{enumerate}
This is why the ceiling rule in \cref{sec:method} requires a generation after the prefill.

\subsection{Host RAM, not VRAM, caused the only outage}\label{sec:hostram}

\begin{figure}[t]
\centering
\begin{tikzpicture}
\begin{axis}[paper,ybar=10pt,height=5.0cm,width=0.72\linewidth,bar width=22pt,
  symbolic x coords={q6,q4},xtick=data,xticklabels={\Qsix{} (\flag{q4\_0} KV),\Qfour{} (\flag{q8\_0} KV)},
  enlarge x limits=0.5,ymin=0,ymax=6600,ylabel={peak host RAM, anonymous (MiB)},
  scaled y ticks=false,yticklabel style={/pgf/number format/1000 sep={,}},
  nodes near coords,nodes near coords style={font=\footnotesize,/pgf/number format/1000 sep={,}},
  legend style={at={(0.5,1.03)},anchor=south,legend columns=2}]
\addplot[fill=cbblue!75,draw=cbblue!60!black] coordinates {(q6,5380) (q4,5396)};
\addplot[fill=cborange,draw=cborange!60!black,postaction={pattern=north east lines,pattern color=black!40}] coordinates {(q6,1752) (q4,1765)};
\legend{MTP $n{=}4$,DFlash2 $n{=}7$}
\end{axis}
\end{tikzpicture}
\caption{Peak host memory of the server process over a 44-turn conversation grown to about 252K
tokens (one soak per configuration, \flag{-ctxcp 4}, one slot). The drafter, not the build, sets the
cost: about 3.6\,GiB between drafters, about 16\,MiB between builds. All four passed.}\label{fig:hostram}
\end{figure}

After the first battery the machine went into daily use serving \Qsix{} with MTP at 262K. During a
long agent session the server's \emph{host} memory grew from 2.9 to 10.0\,GiB while VRAM stayed flat.
Swap filled and prefill fell from 839 to 0.15\tps. The cause: llama.cpp keeps context checkpoints
(snapshots that allow a conversation to roll back without reprocessing, needed because recurrent
layers cannot be rewound \citep{llamacpp_pr20288_checkpoints}) in host RAM, and with MTP each one
includes the draft cache for the whole sequence. Measured from a verbose trace: \textbf{150\,MiB plus
about 3.9\,KiB per token, 1.14\,GiB each at 252K tokens.} I had raised the checkpoint count from 4 to
32 because an A/B test showed slightly faster decoding at identical VRAM. That was true and beside the
point: no VRAM measurement, and no ceiling test made of one prefill and one generation, can see this.

The fix was \flag{-ctxcp 4}, a container memory limit, and a multi-turn soak test as a gate for any
serving change: 44 turns growing to 252K tokens, watching host memory, swap and prefix-cache reuse.
Under that soak (\cref{fig:hostram}) MTP peaks at about 5.4\,GiB of host memory and DFlash2 at about
1.75\,GiB on both builds. On a 14\,GiB host serving a 262K window, the drafter is a host-memory
decision before it is a speed decision. Each additional concurrent slot would carry its own
checkpoints; that was not measured.

\begin{takeaway}
The native 262K window fits on 2$\times$16\,GB, with a 4-bit cache for \Qsix{} and an 8-bit cache for
\Qfour. It takes a hand-tuned layer split that must be re-swept per configuration, a ceiling test that
generates after a deep prefill, and a soak test that watches host RAM.
\end{takeaway}
